\section{Argument soundness}
\label{sec:soundness}

This section gives a concrete classical random-oracle argument-soundness
bound for the relation in Definition~\ref{def:spend-relation}.  The
finite-parameter code reduction and the state-restoration property needed by
the Fiat--Shamir compiler are stated explicitly as premises below.  Section~\ref{sec:knowledge-soundness} states, under an added round-by-round knowledge premise, the extractor that upgrades this to an argument of knowledge and the theft-resistance corollary that follows.

\subsection{Security experiment and metric}

Let \(\Pi_{\mathrm{spend}}\) denote the interactive public-coin protocol obtained by
replacing each Fiat--Shamir challenge with an independent verifier message.
The adversary may choose the public pre-state and statement \(x\), receive the
verifier messages adaptively, and return all oracle messages and openings.  It
wins if the verifier accepts while
\(
  \{w:\RProfile(x,w)=1\}=\varnothing
\).
Account-state rejection can only reduce this event, so the algebraic analysis
may conservatively omit the direct pre-state checks.

\begin{remark}[Statement-encoding change]
The v4 statement encoding extends the statement-digest preimage by the
32-byte deployment domain of Definition~\ref{def:spend-statement}.  The
digest remains one 32-byte SHA-256 value absorbed at the same transcript
position, so the Fiat--Shamir event ledger below is unchanged: no absorb
boundary, challenge, or grinding predicate is added or moved by the new
field.
\end{remark}

For the Fiat--Shamir protocol, let a classical adversary make
\(T\in[1,2^{128}]\) queries to the 256-bit random oracle, including queries
made after it has produced candidate proof bytes.  We report
\begin{equation}
 \operatorname{Adv}^{\mathsf{snd}}_{\mathrm{spend}}(T)
 =\frac{\Pr[\text{the adversary produces an accepting false statement}]}{T},
 \label{eq:work-normalized-soundness}
\end{equation}
the success probability per random-oracle query.  This normalization is
appropriate for a protocol whose explicit grinding work is implemented by
random-oracle search, and it is the metric used for the numerical claim.

\subsection{Code transport and list binding}

Write
\[
 \mathcal C'=
 \{p|_{G'_{19}}:p\in L'_{10}(\mathbb K)\}.
\]
This is the exact 1,024-coefficient subcode evaluated on the
\(2^{19}\)-symbol circle coset.  S-two Theorem~5 and
Equations~(21)--(22) give \(\dim L'_{10}=2^{10}\) and
\(L'_{10}\subseteq L_{10}\).  Its Definition~6, Lemma~1, and Section~1.4
identify the ambient \(L_{10}\) circle code, under the rational circle
parametrization and a nonzero coordinate scaling, with a generalized
Reed--Solomon code of degree at most 1,024; its Theorem~7 supplies the
corresponding Johnson-regime list-decoding result
\cite{CarmonEtAl2026Stwo}.  Coordinate permutation and nonzero diagonal
scaling preserve Hamming distance, agreement sets, and list sizes.  Passing
to \(\mathcal C'\) can only shrink a decoding list.

\begin{lemma}[Circle-to-GRS transport]
\label{lem:circle-grs-transport}
For the \(2^{10}\)-coefficient, \(2^{19}\)-symbol \Profile{} circle code,
the implemented circle-to-line map identifies each relevant affine code
slice with a scaled GRS slice having the same relative distance and the same
agreement sets.  The four implemented arity-four fiber folds commute with
this identification.
\end{lemma}

\begin{proof}
The circle parametrization pairs inverse points into the implemented
arity-four fibers.  Restricting a polynomial to the normalized circle and
then applying \(t=2x^2-1\) gives the line evaluation used by the verifier.
On the selected coset the 131,072 fiber roots are distinct and none is one;
the release evaluator checks all of them.  Evaluation therefore differs
from the corresponding GRS evaluation only by the fixed order and nonzero
normalizers.  Expanding the implemented fold gives the degree-three powers
generator \(a_0+\alpha a_1+\alpha^2a_2+\alpha^3a_3\); the coordinate
divisions are the invertible local FFT change of basis.  S-two Lemma~4 and
WHIR Lemma~4.13/Theorem~4.20 then give fold/list commutation for this powers
generator, while the terminal constraint transports the linear subcode
condition~\cite{CarmonEtAl2026Stwo,ArnonChiesaFenziYogev2025WHIR}.
\end{proof}

The soundness ledger uses proven Johnson list decoding and the
polynomial-generator MCA regime, rather than a capacity conjecture.  The
required multicodeword-agreement bounds are the one-domain specialization of
S-two Protocol~3, Theorem~19, and Lemmas~3--4.  For the scalar-powers
generator, the corresponding Reed--Solomon MCA bound is Bordage et al.,
Definition~9.1 and Lemma~9.3; their Theorem~9.2 extends the result to every
polynomial generator
\cite{CarmonEtAl2026Stwo,BordageChiesaGuanManzur2025,Haboeck2025MCA}.

\begin{lemma}[Johnson/MCA batching]
\label{lem:johnson-mca-batch}
Suppose the polynomial-batching, OOD, and local algebraic bad events listed in
Table~\ref{tab:soundness-events} do not occur.  Then the two layer-zero
commitments and their \(\gamma\)-batch in
Equation~\eqref{eq:spend-batch} determine one member of the transported
rate-\(1/512\) code list, and every accepted fold and terminal value belongs
to the corresponding folded list.
\end{lemma}

\begin{proof}
The width-29 scalar-power batch is a polynomial-generator family.  Outside
its collision event, the MCA theorem reduces simultaneous proximity of the
committed lanes to proximity of their single \(\gamma\)-combination.  The
rate-\(1/512\) parameters lie inside the published proven-Johnson regime.  By
Lemma~\ref{lem:circle-grs-transport}, the bound transfers to the implemented
circle code.  Linear arity-four folding carries every candidate list element
to the next list.  The two OOD samples and relation sumcheck at each layer
bind the claimed list member; their exceptional events are charged
separately in the ledger.  Finally, the four-coefficient terminal check fixes
the last folded polynomial.
\end{proof}

\begin{lemma}[Local algebraic binding]
\label{lem:local-algebraic-binding}
Conditioned on the code/list conclusion of
Lemma~\ref{lem:johnson-mca-batch}, an accepting transcript either encodes a
trace satisfying Proposition~\ref{prop:relation-realization}, or one of the
local events in rows 5--14 of Table~\ref{tab:soundness-events} occurs.
\end{lemma}

\begin{proof}
The relation OOD mixers bind the batched semantic, copy, lookup, and mask
identities.  A nonzero \(\gamma\) separates the 29 generator lanes and the
three statement points.  Tagged-tuple compression, copy denominators, and the
\(\theta\)-batch are sound outside their displayed collision or pole events.
The zerocheck equality point and ten degree-27 rounds bind the trace
constraints, while the \(H_1\) and nonzero-\(\eta\) rows bind the masked
sumcheck.  When these events are absent, all constraints and public terminal
equalities hold, so Proposition~\ref{prop:relation-realization} supplies a
valid witness.
\end{proof}

The preceding proofs identify the intended reduction, but the generated
calculator and regression fixtures are not a formal verification of every
imported-theorem hypothesis.  The concrete theorem therefore isolates the
finite-parameter step rather than silently treating executable agreement as
a proof.

\begin{assumption}[Pinned finite-parameter reduction]
\label{ass:pinned-finite-reduction}
For the exact code, transcript, and verifier serialized by the frozen
artifact, Lemmas~\ref{lem:circle-grs-transport}--
\ref{lem:local-algebraic-binding} hold with the event partition and bounds in
Table~\ref{tab:soundness-events}.  In particular, the premise covers the
rate-\(1/512\) Johnson-list and polynomial-generator MCA hypotheses, the two
OOD samples, all four folds and their list sizes, the terminal condition, and
the assignment of every semantic, copy, lookup, range, mask, and sumcheck
constraint to rows 5--14 of the table.  The premise is uniform over every
false public statement and every reachable public transcript prefix.
\end{assumption}

\subsection{Selector scope}

The selector byte is absorbed only after both commitments, all OOD data,
folds, the terminal polynomial, and the final work nonce have been fixed.
Let \(B\) denote the union of all bad events determined before that byte, and
let \(M_b\) be the final query-miss event for branch \(b\).

\begin{lemma}[Three-branch selector containment]
\label{lem:selector-soundness}
For every possibly adversarial function that selects
\(b\in\{0,1,2\}\) after observing all three schedules,
\[
 B\lor M_b\quad\Longrightarrow\quad B\lor M_0\lor M_1\lor M_2.
\]
Consequently the three-branch construction changes only the final query-miss
term by a factor of at most three; it does not replicate any common-prefix
term.
\end{lemma}

\begin{proof}
The implication is pointwise and does not require independence between the
selector and the schedules.  The common-prefix transcript replay additionally
checks equality of the prefix challenges, OOD points, mixers, fold challenges,
and pre-final-grinding state across all branches.  A union bound over
\(M_0,M_1,M_2\) completes the argument.
\end{proof}

The release theorem below takes the still more conservative course of
multiplying the complete transformed ledger by three.  This avoids relying on
the tighter scope allocation in its headline value.

\subsection{Interactive event ledger}

Table~\ref{tab:soundness-events} lists bounds \(2^{-b_i}\) for the 16
separately named failure families.  The final-query row shows both the
unselected one-branch value used to form the conservative ledger and, in
parentheses, its factor-three selector-local value.  The four fold rows have
already been unioned into one table entry.

For the query row, take Johnson multiplicity \(m=10\),
\(\rho=1/512\), and
\[
 \alpha=\left(1+\frac{1}{2m}\right)\sqrt{\rho}.
\]
Among \(N=131072\) fibers, the agreement cap is
\(A=\lfloor\alpha N\rfloor=6082\).  Sampling 18 distinct fibers therefore
has miss probability at most
\[
 \prod_{i=0}^{17}\frac{6082-i}{131072-i}.
\]
The 32-bit final work predicate raises the corresponding one-branch row to
111.7687016344 bits.

\begin{table}[t]
\centering
\small
\caption{\Profile{} interactive soundness ledger.  Entries are
\(-\log_2\) upper bounds on the named event or conservative policy reserve.}
\label{tab:soundness-events}
\begin{tabular}{@{}lr@{}}
\toprule
event & bits \\
\midrule
width-29 polynomial batch & 107.3160201144 \\
four-fold union & 108.9854322658 \\
final q18 miss, one branch (three-branch local union) & 111.7687016344 (110.1837391336) \\
two-sample OOD-list union & 213.1000183949 \\
24 relation OOD mixers & 119.4150374966 \\
nonzero \(\gamma\) and three-point batch & 119.0931094017 \\
inactive-copy \(\gamma\) claim & 119.1926450753 \\
atomic tuple compression & 112.3968373178 \\
atomic copy/range poles & 111.7627900366 \\
atomic \(\theta\) collision & 119.4150374966 \\
zerocheck equality point & 120.6780719024 \\
zero-sum \(H_1\) helper & 123.9999999973 \\
nonzero \(\eta\) for the original mask & 123.9999999973 \\
ten degree-27 zerocheck rounds & 115.9231844003 \\
Poseidon2 semantic-binding policy reserve & 124 \\
SHA-256 instantiation policy reserve & 128 \\
\bottomrule
\end{tabular}
\end{table}

Rows 1--14 are the algebraic-event partition in
Assumption~\ref{ass:pinned-finite-reduction}.  Rows 15--16 are release-policy
slack, not collision probabilities derived for an adversary making
\(2^{128}\) queries.  Argument soundness for the exact relation treats the
Poseidon2 gates as part of \(\RProfile\), while the theorem below treats
\(\mathcal H\) as an ideal random oracle.  Omitting these two positive reserve
terms would only strengthen the numerical bound.  Collision resistance of
Poseidon2 is instead needed when a note commitment is interpreted as a
binding name for application data outside the exact relation.

Using the one-branch final-miss row and unioning all 16 events gives
\begin{equation}
 \varepsilon_{\rm round}^{(0)}
   =\sum_i 2^{-b_i}
   \le 2^{-106.79080600295417}.
 \label{eq:coarse-interactive-error}
\end{equation}
Under Assumption~\ref{ass:pinned-finite-reduction}, the complement of this
union and Lemmas~\ref{lem:johnson-mca-batch}
and~\ref{lem:local-algebraic-binding} imply that \(\RProfile\) is satisfiable.

\subsection{Fiat--Shamir instantiation}

The BCS compiler is characterized by soundness against state-restoration
attacks, which is stronger than ordinary interactive soundness
\cite{BenSassonChiesaSpooner2016,BlockGarretaKatzThalerTiwariZajac2023}.
A state-restoration adversary may save the verifier and oracle state at a
public-message boundary, restore a saved state, and request a fresh
continuation; its win event is acceptance of a false statement on any
resulting branch.  Let \(\varepsilon_{\rm sr}\) denote the maximum win
probability in the exact 32-boundary experiment for the serialized
transcript.

\begin{assumption}[Boundary state restoration]
\label{ass:boundary-restoration}
For every false statement, every reachable saved state at each of the 32
boundaries, and every classical state-restoration adversary, the conditional
escape events are covered by the event partition in
Assumption~\ref{ass:pinned-finite-reduction}.  Consequently
\[
 \varepsilon_{\rm sr}\le\varepsilon_{\rm round}^{(0)}.
\]
This is a round-by-round premise; it is not inferred from standard
interactive soundness.
\end{assumption}

We use the BCS Fiat--Shamir analysis for IOPs
\cite{BenSassonChiesaSpooner2016}, in the explicit bounded-query form of
S-two Theorem~22~\cite{CarmonEtAl2026Stwo}.  Under
Assumption~\ref{ass:boundary-restoration}, its instantiation with \(R=32\)
message boundaries and \(\lambda=256\) gives
\begin{equation}
 \varepsilon_{\mathrm{BCS}}(T)
 \le (T+R)\varepsilon_{\rm round}^{(0)}
       +\frac{3(T^2+1)}{2^\lambda}.
 \label{eq:bcs-raw}
\end{equation}
At the upper query budget, the right-hand side of
Equation~\eqref{eq:bcs-raw} exceeds one and is therefore vacuous as a bound
on raw success probability.  The released claim is only the explicitly
defined work-normalized metric in
Equation~\eqref{eq:work-normalized-soundness}; it does not book the raw bound
as a security floor.
Dividing by \(T\) gives the metric in
Equation~\eqref{eq:work-normalized-soundness}:
\begin{equation}
 b(T):=\frac{\varepsilon_{\mathrm{BCS}}(T)}{T}
 \le
 \left(1+\frac{32}{T}\right)\varepsilon_{\rm round}^{(0)}
 +\frac{3(T+1/T)}{2^{256}}.
 \label{eq:bcs-normalized}
\end{equation}
At \(T=1\), the multiplier of the round error is therefore 33.  The
right-hand side is convex for positive \(T\), so its maximum on
\([1,2^{128}]\) is attained at an endpoint.  Direct evaluation yields
\[
 -\log_2 b(1)=101.74641188359571,
 \qquad
 -\log_2 b(2^{128})=106.79080421773332.
\]

\begin{theorem}[Work-normalized classical-ROM soundness]
\label{thm:bcs-soundness}
Assume the pinned finite-parameter reduction
(Assumption~\ref{ass:pinned-finite-reduction}), the 32-boundary
state-restoration property (Assumption~\ref{ass:boundary-restoration}), and
SHA-256 modeled as a 256-bit random oracle.  For every classical adversary making
\(1\le T\le2^{128}\) oracle queries,
\[
 \operatorname{Adv}^{\mathsf{snd}}_{\mathrm{spend}}(T)
 \le 3b(T)\le 2^{-100.16144938287455}.
\]
Thus \Profile{} has a conservative 100.161-bit argument-soundness floor in
the work-normalized metric.
\end{theorem}

The work-normalized floor is a per-query bound.  The cumulative raw advantage
at a fixed query budget is \(3\varepsilon_{\mathrm{BCS}}(T)\); it grows with
\(T\) and becomes vacuous as \(T\) approaches the round error, as every
grinding-based argument does.  Table~\ref{tab:raw-advantage} reports both so
the metric is not misread as a raw \(2^{-100}\) forgery probability at every
budget.

\begin{table}[t]
  \centering
  \small
  \caption{Raw false-acceptance advantage \(3\varepsilon_{\mathrm{BCS}}(T)\)
  by query budget, against the uniform per-query floor.}
  \label{tab:raw-advantage}
  \begin{tabular}{r r}
    \toprule
    Query budget \(T\) & \(-\log_2\) raw advantage \\
    \midrule
    \(2^{40}\) & 65.21 \\
    \(2^{64}\) & 41.21 \\
    \(2^{80}\) & 25.21 \\
    \(2^{100}\) & 5.21 \\
    \(\gtrsim 2^{105.2}\) & vacuous \\
    \midrule
    per-query floor (all \(T\)) & 100.161 \\
    \bottomrule
  \end{tabular}
\end{table}

\begin{proof}
Assumptions~\ref{ass:pinned-finite-reduction} and
\ref{ass:boundary-restoration}, together with
Equation~\eqref{eq:coarse-interactive-error}, bound the state-restoration
experiment.  Equation~\eqref{eq:bcs-raw} then applies the 32-boundary
Fiat--Shamir transform.
Multiplication by three is a whole-ledger upper bound for the three possible
post-final query schedules.  Evaluating the two interval endpoints after this
factor gives 100.16144938287455 and 105.20584171701216 bits, respectively;
the former is the minimum.  Every accepted false statement is included in
this event by the preceding lemmas.
\end{proof}

\subsection{Knowledge soundness and theft resistance}
\label{sec:knowledge-soundness}

Theorem~\ref{thm:bcs-soundness} is a statement about the \emph{existence} of a
satisfying witness.  For a spend, the property a pool needs is stronger: only a
party that knows a note's authorization secret should be able to produce an
accepting transition that consumes it.  We record the additional premise under
which the argument becomes an argument of knowledge, and the theft-resistance
statement that follows from it and from the relation's binding structure.  Like
the soundness premises, the knowledge premise is assumed, not proved.

\begin{assumption}[State-restoration knowledge extraction]
\label{ass:sr-knowledge}
There is an expected-polynomial-time extractor \(\mathsf{Ext}\) with the
following property.  For the exact code, transcript, and verifier of the frozen
artifact, and for every classical state-restoration prover \(P^\ast\) that from
any reachable saved state at each of the 32 boundaries drives the interactive
protocol to an accepting terminal with probability exceeding
\(\varepsilon_{\rm round}^{(0)}\), \(\mathsf{Ext}\), given \(P^\ast\) together
with its transcript of random-oracle queries and answers, outputs codeword
openings and constraint assignments that satisfy every pinned event of
Table~\ref{tab:soundness-events}, except with probability at most
\(\varepsilon_{\rm round}^{(0)}\).  This is the round-by-round knowledge
analogue of Assumption~\ref{ass:boundary-restoration}; it is not inferred from
Theorem~\ref{thm:bcs-soundness}.
\end{assumption}

\begin{theorem}[Conditional non-interactive argument of knowledge]
\label{thm:argument-of-knowledge}
Assume Assumptions~\ref{ass:pinned-finite-reduction},
\ref{ass:boundary-restoration}, and~\ref{ass:sr-knowledge}, and SHA-256 as a
programmable random oracle.  There is a straight-line extractor \(E\) such that
for every prover that outputs an accepting proof \(\pi\) for a statement \(x\)
while making \(T\le2^{128}\) oracle queries, \(E\), given that prover's
query--answer transcript, outputs \(w\) with \(\RProfile(x,w)=1\), except with
probability at most \(3b(T)\le2^{-100.161}\).
\end{theorem}

\begin{proof}
The BCS Fiat--Shamir transform of a state-restoration interactive oracle proof
is an argument of knowledge in the random-oracle model
\citep{BenSassonChiesaSpooner2016}: the extractor recovers the Merkle-committed
oracle messages from the prover's query transcript, and by
Assumption~\ref{ass:sr-knowledge} their openings at the sampled positions
determine codeword openings and assignments satisfying the pinned constraints.
The pinned encoding of Section~\ref{sec:construction} maps that interactive
witness to a relation witness \(w\) for \(\RProfile\).  Extraction is
straight-line from the query transcript, so it uses no rewinding and its
failure is the same state-restoration escape event bounded in
Theorem~\ref{thm:bcs-soundness}; the \(3b(T)\) floor therefore carries
unchanged.
\end{proof}

\begin{corollary}[Authorization possession and theft resistance]
\label{cor:theft-resistance}
Fix a note whose input commitment lies in the committed tree, whose owner key
is \(pk_{\mathrm{in}}=H_{\mathsf{own}}(k_\nu)\) for the holder's secret
\(k_\nu\) by Equation~\eqref{eq:owner-key}, and whose public nullifier is
\(\nu=H_{\mathsf{nul}}(k_\nu,r_{\mathrm{in}})\) by
Equation~\eqref{eq:nullifier}.  Let \(A\) be any probabilistic
polynomial-time party that outputs an accepting spend consuming that note --- an
accepting proof for a statement carrying \(\nu\) --- while making \(T\le2^{128}\)
oracle queries.  Then, except with probability at most
\(3b(T)+\varepsilon_{\rm cr}\), the extractor of
Theorem~\ref{thm:argument-of-knowledge} yields a witness containing a value
\(k\) with \(H_{\mathsf{own}}(k)=pk_{\mathrm{in}}\) and
\(H_{\mathsf{nul}}(k,r_{\mathrm{in}})=\nu\); by collision resistance of
Poseidon2, with reserve \(\varepsilon_{\rm cr}\) in the soundness ledger,
\(k=k_\nu\).  Hence any efficient party that produces an accepting spend of the
note knows its authorization secret, and no efficient party lacking \(k_\nu\)
produces an accepting spend except with probability at most
\(3b(T)+\varepsilon_{\rm cr}\).
\end{corollary}

\begin{remark}[Deployed pools and scope]
\label{rem:sim-extractability}
A deployed pool additionally exposes honest spends to the adversary.  The
witness-free simulator of Theorem~\ref{thm:real-vs-sim23} programs the oracle
only at fresh inputs, while the extractor reads the adversary's own query
transcript, so simulation and extraction do not interfere except with
negligible probability; the argument is therefore simulation-extractable and
Corollary~\ref{cor:theft-resistance} holds against an adversary with oracle
access to honest spends.  The guarantee is conditional on
Assumption~\ref{ass:sr-knowledge}, the knowledge analogue of the
state-restoration premise, which is assumed rather than proved, and it concerns
the production of accepting proofs only: it does not address key management,
secret-key leakage, or the local and network observables excluded from the
declared view of Definition~\ref{def:spend-complete-view}.
\end{remark}
